% Part: first-order-logic
% Chapter: beyond
% Section: modal-logics

\documentclass[../../../include/open-logic-section]{subfiles}

\begin{document}

\olfileid{fol}{byd}{mod}

\olsection{Modale logica}

Beschouw het volgende voorbeeld van een voorwaardelijke zin:
\begin{quote}
  Als Jeremy alleen in die kamer is, dan is hij dronken en naakt en
  danst hij op de stoelen.
\end{quote}
Dit is een voorbeeld van een voorwaardelijke bewering die materieel waar kan
zijn en niettemin misleidend kan zijn. De zin lijkt immers te suggereren dat er
een sterker verband tussen het antecedens en de conclusie bestaat dan enkel dat
het antecedens onwaar of het consequens waar is. De formulering suggereert met
andere woorden dat de bewering niet alleen in deze specifieke wereld waar is
(waar zij triviaal waar kan zijn doordat Jeremy niet alleen in de kamer is),
maar bovendien dat de conclusie waar \emph{zou zijn geweest} \emph{als} het
antecedens waar was geweest. Men kan de bewering dus zo opvatten dat zij niet
alleen in deze wereld, maar in elke ``mogelijke'' wereld waar is; ofwel dat zij
\emph{noodzakelijk} waar is, in plaats van slechts waar in deze specifieke
wereld.

De modale logica is ontwikkeld om dit soort noodzakelijkheid te duiden. De
modale propositielogica ontstaat door aan de gewone propositielogica een
boxoperator toe te voegen; dat wil zeggen: als $!A$ !!a{formula} is, dan is ook
$\Box !A$ dat. Intuïtief beweert $\Box !A$ dat $!A$ \emph{noodzakelijk} waar is,
of waar in elke mogelijke wereld. $\Diamond !A$ wordt gewoonlijk opgevat als
een afkorting van $\lnot \Box \lnot !A$ en kan worden gelezen als de bewering
dat $!A$ \emph{mogelijk} waar is. Uiteraard kan modaliteit ook aan de
predicatenlogica worden toegevoegd.

Kripke-!!{structure}s kunnen als semantiek voor de modale logica dienen; Kripke
ontwierp deze semantiek oorspronkelijk zelfs met de modale logica voor ogen. In
plaats van ons tot partiële ordeningen te beperken, nemen we algemener een
verzameling ``mogelijke werelden'', $P$, en een binaire
``toegankelijkheidsrelatie'' $\Atom{R}{x,y}$ tussen werelden. Intuïtief beweert
$\Atom{R}{p,q}$ dat de wereld~$q$ verenigbaar is met~$p$: als we ons in
wereld~$p$ ``bevinden'', moeten we rekening houden met de mogelijkheid dat de
wereld als~$q$ had kunnen zijn.

Modale logica wordt soms een ``intensionele'' logica genoemd, tegenover een
``extensionele'' logica. De beoogde semantiek voor een extensionele logica,
zoals de klassieke logica, verwijst slechts naar één wereld, de ``werkelijke'';
de semantiek voor een intensionele logica berust daarentegen op een uitgebreidere
ontologie. Modaliteit kan niet alleen noodzakelijkheid !!{structure}ren, maar ook
andere talige constructies, waarbij $\Box$ en $\Diamond$ naargelang de toepassing
anders worden geïnterpreteerd. Bijvoorbeeld:
\begin{enumerate}
\item In bewijsbaarheidslogica wordt $\Box !A$ gelezen als ``$!A$ is bewijsbaar''
  en $\Diamond !A$ als ``$!A$ is consistent.''
\item In epistemische logica kan men $\Box !A$ lezen als ``ik weet
  $!A$'' of ``ik geloof $!A$.''
\item In tijdslogica kan men $\Box !A$ lezen als ``$!A$ is altijd
  waar'' en $\Diamond !A$ als ``$!A$ is soms waar.''
\end{enumerate}

We willen de logica uitbreiden met regels en axioma's voor modaliteit. Het
systeem $\Log{S4}$ bestaat bijvoorbeeld uit de gewone axioma's en regels van de
propositielogica, samen met de volgende axioma's:
\begin{align*}
& \Box (!A \lif !B) \lif (\Box !A \lif \Box !B)\\
& \Box !A \lif !A \\
& \Box !A \lif \Box \Box !A
\intertext{en de regel ``concludeer uit $!A$ de formule $\Box !A$.''
  $\Log{S5}$ voegt daaraan het volgende axioma toe:}
& \Diamond !A \lif \Box \Diamond !A
\end{align*}
Varianten van deze axioma's kunnen voor verschillende toepassingen geschikt
zijn; S5 wordt bijvoorbeeld gewoonlijk geacht het begrip logische
noodzakelijkheid te karakteriseren. Het aantrekkelijke is dat doorgaans een
semantiek kan worden gevonden waarvoor het !!{derivation} systeem correct en
volledig is, door de toegankelijkheidsrelatie in de Kripke-!!{structure}s op
natuurlijke wijze te beperken. Zo correspondeert $\Log{S4}$ met de klasse van
Kripke-!!{structure}s waarin de toegankelijkheidsrelatie reflexief en transitief
is. $\Log{S5}$ correspondeert met de klasse van Kripke-!!{structure}s waarin de
toegankelijkheidsrelatie {\em universeel} is, dat wil zeggen dat elke wereld
vanuit elke andere wereld toegankelijk is; $\Box !A$ geldt dus dan en slechts
dan als $!A$ in elke wereld geldt.

\end{document}
